\chapter{作品概述}
\section{项目背景}
随着信息技术的发展，数据传输面临的风险日趋复杂。
仅保护传输内容已不足以应对这些风险，传输行为本身也可能需要避免引起第三方不必要的关注。
因此，现代数据传输提出了双重需求：一是安全性，即防止未授权方获取受保护信息；
二是隐蔽性，即减少可被观测的痕迹，避免暴露敏感信息正在传输~\cite{1-wang-moulin-2008}。
传统加密技术具有强保密性，并被广泛用于保护消息内容。
然而，仅有保密性并不能隐藏受保护传输的存在。
图像隐写通过在自然载体中隐藏信息，满足了这一补充性需求[1][2]。
将二者直接结合虽可同时提供安全性与隐蔽性，但通常需要额外组件，
并带来准备和部署上的负担[3][4]。因此，如何直接复用现有计算资源，构建安全且隐蔽的传输机制，成为一个值得探索的方向。

深度神经网络（DNN）已广泛部署于各类信息系统[5]，
并日益成为重要计算基础设施的一部分。
同时，它们在密码学[6][7]和图像隐写[8]中的应用也表明，神经计算在信息保护方面具有潜力。
常见视觉DNN通过非线性变换将输入图像映射为高维神经元状态，为信息表示提供了丰富空间[9]。
这引出一个值得探究的问题：现有DNN能否利用其神经元状态表示敏感信息，从而实现由密钥控制的恢复，
并实现类图像的隐蔽传输？受此问题启发，我们探索了一种基于表示的框架。该框架不同于传统密码学与隐写学范式，
旨在同时提供安全性与隐蔽性。它复用已部署于各类信息系统中的现有DNN，
而无需修改其参数或功能。具体而言，我们将消息表示为选定神经元的目标状态，
并通过对抗优化将随机图像迭代修改为外观自然的密文图像。
当该密文图像与预共享图像密钥结合时，可在DNN中诱导出目标状态，
从而实现密钥控制的消息恢复，同时保持类图像的传输隐蔽性。

然而，设计此类框架面临两个主要挑战。第一，难以将DNN精确引导至目标状态。
尽管对抗扰动能够改变DNN状态[10][11]，但神经元之间复杂的依赖与耦合关系，
使得将输入修改精确转化为DNN状态变化十分困难。
一些神经元可能对这类修改不敏感，无法达到目标值[12][13]，我们称其为不响应神经元。
第二，生成的密文图像可能泄露额外信息。其中包含用于引导目标DNN状态的细微图像修改。
然而，即使没有图像密钥，这些修改仍可能通过视觉变化[14]或DNN状态偏差泄露信息。
因此，如何在利用图像密钥诱导目标状态的同时抑制这种泄露，是另一项挑战。

为解决上述挑战，我们提出了 SSIR，一种由图像密钥控制 DNN 状态的安全隐蔽信息表示框架。
在该框架中，发送方与接收方预先秘密共享一个图像密钥，而 DNN 与编码方法均公开且确定。
发送方将随机图像与图像密钥叠加，得到融合图像；将其输入 DNN，并迭代修改该融合图像，
直至 DNN 呈现表示消息的目标状态。随后，发送方从融合图像中减去图像密钥，得到密文图像，
并通过公共信道传输。接收方将密文图像与图像密钥叠加，得到融合图像；
将其输入 DNN，并通过解码所得 DNN 状态恢复消息。

\section{技术背景}


\section{相关工作}
在神经网络与密码学的交叉研究中，已有工作大致可分为两类：一类探索基于学习的信息保护机制，另一类研究密码操作的神经实现。
Abadi 和 Andersen [6] 提出了对抗神经密码学，通过联合训练负责保护与恢复的合作网络，
并与敌对网络对抗。该方法利用学习得到的变换实现信息保护，其中专门的网络训练是框架构建的一部分。
Gérault 等人 [7] 则研究了如何在深度神经网络中安全实现密码操作。
他们指出，当输入由二进制值扩展到任意实数时，
直接用神经网络实现的分组密码可能变得脆弱。
为弥合离散密码计算与连续神经计算之间的差异，
他们提出了基于 ReLU 的实现，并保证了正确性与安全性。
该工作表明将密码功能融入神经网络具有可行性，同时凸显了两个计算领域之间接口设计的重要性。

基于图像的信息传输是另一个相关方向 [2], [16]–[18]。
HiDDeN [3] 通过联合训练嵌入网络与提取网络，将消息隐藏于图像载体中。
固定神经网络隐写 [19] 则通过固定网络优化图像以恢复消息，从而减少对专用网络训练的依赖。
Ghamizi 等人 [20] 利用对抗扰动实现隐写通信，Luo 等人 [8] 在固定网络隐写中引入密钥控制的恢复。
这些工作从不同角度探索了图像优化、神经响应与恢复控制的组合。

与上述工作不同，SSIR 利用图像密钥控制现有公开分类 DNN 的状态，
实现安全且隐蔽的信息表示。它将图像密钥融合与分组编码相结合，
以提高状态诱导的可靠性，同时保持所复用 DNN 的参数与原始功能不变。
联合优化可抑制视觉泄漏与未加密 DNN 状态痕迹，形式化分析则刻画了基于查询的区分优势。
该设计将可靠状态诱导、泄漏抑制与安全性刻画统一在一个信息表示框架中。
\section{作品特色}
通过分析已有方案的优缺点，我们提出了 SSIR——一种由图像密钥控制现有公开 DNN 状态的安全隐蔽信息表示框架。
该框架基于图像密钥融合、分组编码和对抗优化实现，具有以下特点：

\begin{itemize}
    \item \textbf{新型信息表示框架}：将机密信息表示为由图像密钥控制的 DNN 状态，
    以实现安全且隐蔽的传输。为精确诱导目标状态，
    本作品将信息编码为目标比特串，并利用分组神经元响应与目标比特之间的可微损失，
    指导对扰动图像的迭代修改，从而生成密文图像。分组编码对相邻神经元响应取平均，
    以缓解不响应神经元的影响，提高信息表示的可靠性。

    \item \textbf{双重泄漏抑制机制}：分析了对手在没有图像密钥时可能从密文图像中获取的额外信息，
    并将其分为两类：直接视觉泄漏，即密文图像与原始图像之间的视觉差异；
    间接 DNN 状态泄漏，即仅由密文图像引起的 DNN 状态变化。
    为减轻上述泄漏，本作品在生成密文图像时设计了基于损失的优化策略，同时抑制视觉痕迹和未加密 DNN 状态偏差。

    \item \textbf{形式化安全分析}：对 SSIR 进行了形式化安全分析，
    分析了输入依赖的 DNN 响应如何支持密钥控制的表示，并刻画了概率密文图像生成与多语义特性。
    通过混合论证，本作品推导了查询相关的 IND-CPA 和 IND-CCA1 优势界，
    并给出了实现有限查询安全所需的分布条件。由此得到的可接受查询限制，
    为配置每个密钥的使用与更新策略提供了基础。

    \item \textbf{高效兼容的实验验证}：在 ImageNet~[15] 和真实照片上评估了 SSIR。
    结果表明，SSIR 能够抵抗多种攻击，对位级失真具有容忍度，
    并通过分组编码提升了生成的可靠性与效率。直接泄漏最高降低 99.79\%，
    间接泄漏平均降低 56.02\%。同时，SSIR 具有较高效率与良好兼容性：
    生成一条 150 位信息的密文图像仅需 300~ms，消息恢复仅需 10~ms，且无需额外模型准备。
    并且可支持多种现有 DNN，而无需修改其参数或原有功能。
\end{itemize}
\section{应用前景分析}
当前，敏感信息传输面临双重困境：传统加密技术虽然能够保护消息内容，但无法隐藏通信行为本身，
密文在公开信道上的频繁出现容易引起第三方注意，图像隐写技术虽然能够把消息隐藏到自然载体中，
却往往依赖专用网络训练或额外组件，部署成本较高，密钥控制、状态诱导可靠性和抗攻击能力也难以兼顾。
随着即时通信、云会议、物联网、应急指挥等场景的发展，会话密钥、控制指令、身份凭证等短小敏感信息、
需要在公开信道上频繁、快速、隐蔽地传递。仅靠加密难以避免通信行为被监测，
仅靠隐写又难以保证安全性与可验证性。因此，能够复用现有计算资源、
同时兼顾保密性与传输隐蔽性的轻量级信息表示机制，将成为未来敏感通信的重要方向。

SSIR 正是面向这一需求提出的。它利用图像密钥控制现有公开 DNN 的神经元状态，
将消息编码为目标状态，并通过对抗优化生成外观自然的密文图像。
接收方只需将同一图像密钥叠加到密文图像上，将合成图像输入 DNN，即可通过解码神经状态恢复消息。

在具体应用上，SSIR 可用于：即时通信和云会议中的会话密钥传递与刷新，降低密钥明文暴露风险；
物联网和边缘设备间的机密控制指令传输，在公开信道上减少被监测和干扰的可能；
应急通信、企业内部敏感指令和身份凭证的隐蔽分发，以及合法合规的秘密投递通信。
对于满足分布条件的实例，部署方可按目标优势预算设置每密钥的生成与恢复上限，
并在同一密钥的会话中累计计数；达到上限前，各方独立生成并秘密预共享替换密钥，
同时将外部恢复访问限制在已分析的威胁模型内。
这符合既定密钥管理指南[26]，也使系统在长期运行中保持可配置的安全边界。

因此，SSIR 不仅是一种信息表示方法，也具备实际部署价值。
它既能服务于企业内部的密钥管理和隐蔽通信，也能用于应急、物联网等对安全性和隐蔽性要求较高的场景，
具有较好的应用前景。

